\section{Evaluation Infrastructure: sin retroalimentación no hay iteración de la alineación abierta}

La sección anterior presentó una gran cantidad de modelos, pero sin un sistema de referencia común. Entre mayo y julio de 2023 aparecieron, en cuestión de semanas, ChatBotArena, AlpacaEval, MT-Bench y Open LLM Leaderboard. La causa no fue que la teoría de la evaluación madurara de repente, sino que los equipos abiertos no podían, como las grandes empresas, contratar personas de forma continua para comparar respuestas. Por eso la evaluación se convirtió en la retroalimentación sustituta del loop de entrenamiento.

\subsection{Por qué surgieron al mismo tiempo cuatro sistemas de evaluación}

Esta sección revisa primero la cronología y la división de funciones. Arena recopila pairwise preference de personas reales: su retroalimentación es lenta, pero tiene alta validez externa. AlpacaEval y MT-Bench usan un LLM judge potente para acelerar la iteración. Leaderboard ejecuta automáticamente tareas estáticas y es el más fácil de integrar en un pipeline de ingeniería. Los cuatro optimizan objetivos distintos y ninguna puntuación única puede sustituirlos.

\lecturefigure{slide-049.jpg}{En 2023, cuatro de las principales aligned-model evaluations abiertas aparecieron casi al mismo tiempo}{official deck, p.~49}

\teachervoice{00:32:14--00:32:48, Nathan explica que estas herramientas surgieron en poco tiempo porque los builders abiertos desperately needed signal, pero no podían pagar de forma continua el costo de las comparaciones humanas como lo hacen las empresas de código cerrado.}

\subsection{ChatBotArena: lenta, pero cercana a la preferencia real de los usuarios}

La subsección anterior colocó los cuatro sistemas de evaluación en un mismo mapa; esta examina primero el que más se acerca a la retroalimentación real de un producto. Arena muestra dos modelos anónimos lado a lado al usuario, recopila victorias, derrotas o empates y después estima su fuerza relativa mediante pairwise ranking. Al leer el diagrama de flujo, conviene seguir primero de dónde viene el prompt, cómo entran los resultados de la comparación en el ranking y cuántos enfrentamientos necesita un modelo nuevo para estabilizarse. Su ventaja es que los prompts vienen de usuarios reales y el judge es una persona; su debilidad es que tanto el muestreo como la convergencia son lentos, de modo que la retroalimentación para el desarrollo puede tardar varios días.

\lecturefigure{slide-050.jpg}{ChatBotArena: comparación anónima entre dos modelos y recopilación de la preferencia real de los usuarios}{official deck, p.~50}

La probabilidad esperada de victoria en Elo puede escribirse como
\[
P(A\succ B)=\frac{1}{1+10^{(R_B-R_A)/400}},
\]
donde $R_A,R_B$ son las calificaciones de los modelos y $P(A\succ B)$ es la probabilidad de que A venza a B. La calificación es una magnitud relativa que depende del grafo de oponentes, de la distribución de usuarios y de la ventana temporal; no conviene comparar directamente el Elo absoluto entre distintas versiones de Arena.

\teachervoice{00:33:14--00:33:56, el ponente considera que Arena es muy importante para los lanzamientos y la estrategia de las empresas, pero demasiado lenta para la ingeniería de entrenamiento: el equipo necesita saber si el modelo mejoró antes de decidir el lanzamiento, por lo que debe complementarse con evaluaciones automáticas más rápidas.}

\subsection{AlpacaEval: preference proxy rápido y judge bias}

AlpacaEval usa prompts fijos, compara la candidate response con la baseline response y después pide a un LLM potente que decida cuál gana. Más muestras producen barras de error más pequeñas; es rápido de ejecutar y fácil de integrar. Sin embargo, la baseline, el judge, la response length y el style modifican la tasa de victorias.

\lecturefigure{slide-051.jpg}{AlpacaEval: LLM preference judging entre la candidate response y la baseline response}{official deck, p.~51}

\lecturefigure{slide-052.jpg}{Ventajas y limitaciones de AlpacaEval: muchas muestras y poco error, pero el judge tiene sesgos}{official deck, p.~52}

Si en $n$ comparaciones independientes la proporción de victorias de la candidate es $\hat p$, el error estándar aproximado es
\[
\operatorname{SE}(\hat p)=\sqrt{\frac{\hat p(1-\hat p)}{n}}.
\]
donde $n$ es el número de comparaciones y $\hat p$ es la win rate observada. Aumentar las muestras reduce el error aleatorio, pero no elimina el sesgo sistemático; por ejemplo, que el judge prefiera respuestas más largas, más corteses o parecidas a su propio estilo.

\lecturefigure{slide-053.jpg}{AlpacaEval 2: cambio a una baseline de GPT-4 e incorporación de length control}{official deck, p.~53}

\begin{warningbox}{Cuatro fuentes de sesgo en LLM-as-a-judge}
\begin{enumerate}
  \item \textbf{Position bias}: la respuesta que aparece primero o al final gana con más facilidad;
  \item \textbf{Length/style bias}: las respuestas extensas y con formato vistoso reciben una preferencia adicional;
  \item \textbf{Self preference}: el judge prefiere salidas parecidas a su propia distribución;
  \item \textbf{Contamination}: el judge puede haber visto el benchmark, la candidate o la reference.
\end{enumerate}
El intercambio aleatorio de posiciones, el control de longitud, el uso de varios judges y la revisión humana por muestreo solo atenúan estos sesgos; no los eliminan por completo.
\end{warningbox}

\teachervoice{00:33:56--00:36:16, Nathan considera AlpacaEval una señal de ingeniería rápida, pero se muestra cauteloso sobre el significado concreto de las puntuaciones de AlpacaEval 2. Cuanto más fácil de optimizar es una evaluación, más necesario es comprobar que el modelo no solo aprendió a complacer al judge.}

\subsection{MT-Bench: cobertura multiturno y fragilidad por muestra pequeña}

AlpacaEval usa pairwise preference para evitar la calibración de puntuaciones absolutas; MT-Bench, en cambio, opta deliberadamente por calificaciones escalares multiturno y por categorías, para observar si el modelo mantiene la calidad después de preguntas de seguimiento. Por eso esta subsección compara dos fuentes de error: la preferencia relativa del pairwise judge y la deriva de una escala de 1 a 10. MT-Bench pide a un LLM judge que califique respuestas de dos turnos, con preguntas que cubren categorías como redacción, razonamiento, matemáticas y programación. Al leer la figura, conviene comprobar primero si el segundo turno depende realmente del primero y después si el promedio por categoría oculta fallos individuales. La estructura multiturno se acerca más al chat que la preference de un solo turno, pero el número de preguntas es pequeño, la escala de calificación del judge es inestable y los equipos pueden ajustar fácilmente sus modelos a las preguntas públicas.

\lecturefigure{slide-054.jpg}{MT-Bench: un modelo potente asigna calificaciones escalares a respuestas multiturno}{official deck, p.~54}

\lecturefigure{slide-055.jpg}{Ventajas, desventajas y riesgo de sobreoptimización de MT-Bench}{official deck, p.~55}

\begin{knowledgebox}{Lectura de la figura: las puntuaciones escalares son más difíciles de calibrar que las pairwise}
Asignar un 7 a una respuesta requiere un criterio absoluto estable, mientras que juzgar si A es mejor que B solo requiere una comparación local. Distintos judges o distintas prompt wording alteran la escala de 1 a 10. Al reportar MT-Bench deben conservarse los resultados por pregunta, la versión del judge y la distribución por categorías, no solo el promedio.
\end{knowledgebox}

\teachervoice{00:36:16--00:37:44, el curso señala que las preguntas de dos turnos de MT-Bench son útiles, pero son pocas, el judge tiene sesgos y, una vez publicadas, pueden optimizarse directamente. Un modelo con puntuación alta puede simplemente ajustarse mejor a esas 80 preguntas.}

\subsection{Open LLM Leaderboard: cómo una herramienta de ingeniería cae bajo la ley de Goodhart}

Leaderboard ejecuta automáticamente varios benchmarks estáticos; en un principio era una herramienta de ingeniería para comparar con facilidad modelos base/fine-tuned. A medida que la clasificación empezó a influir en las descargas y en la reputación, los datos de entrenamiento comenzaron a incluir contenido con el estilo de los benchmarks e incluso preguntas de prueba, y las puntuaciones estáticas fueron perdiendo la capacidad de distinguir la chat quality real.

\lecturefigure{slide-056.jpg}{Open LLM Leaderboard: pipeline automático de benchmarks y clasificación pública}{official deck, p.~56}

\lecturefigure{slide-057.jpg}{Combinación de evaluaciones: Arena es lenta y realista; las eval automáticas son rápidas y se prestan a la manipulación}{official deck, p.~57}

\begin{importantbox}{La forma que toma la ley de Goodhart en la evaluación de modelos}
Cuando una metric pasa de «objetivo de medición» a «objetivo de entrenamiento», el modelo aprende la parte aprovechable de la metric. La solución no es buscar un único benchmark que nunca pueda manipularse, sino construir una combinación de evaluaciones: prompts ocultos o dinámicos, preferencia real de los usuarios, estratificación entre capacidades y seguridad, auditoría humana y reemplazo periódico de las tareas ya saturadas.
\end{importantbox}

\teachervoice{00:37:44--00:39:45, Nathan explica que Leaderboard fue en un principio una herramienta de ingeniería, pero la contaminación y el gaming redujeron su valor; su recomendación de combinación es usar Arena como señal pública de largo plazo y AlpacaEval/MT-Bench como señales rápidas de desarrollo, en lugar de buscar una puntuación universal.}

\begin{lstlisting}[caption={Estructura mínima de registro de un portfolio de evaluación}]
evaluation = {
    "prompt_distribution": "real_users | fixed_set | hidden_set",
    "judge": "human | model_name_and_version",
    "baseline": "model_or_none",
    "aggregation": "elo | win_rate | scalar_mean",
    "latency": "minutes | days",
    "known_biases": ["length", "position", "contamination"]
}
\end{lstlisting}

\subsection{Resumen de la sección}

Los cuatro sistemas de evaluación sirven a escalas de tiempo distintas: Arena respalda los lanzamientos y la comparación pública; AlpacaEval/MT-Bench respaldan la iteración rápida, y Leaderboard respalda la regresión automática de capacidades. Cada herramienta parte de supuestos sobre el prompt, el judge, la baseline y la aggregation. La madurez de la alineación abierta no consiste en encontrar una puntuación perfecta, sino en saber qué puede responder cada puntuación y qué no.
